% SPDX-License-Identifier: Apache-2.0
\section{A Register, Read in Place}
\label{sec:x-inplace}

A register is \code{Copy}, and it has the operators and the
conversions of its value, so it reads in place: \code{self.acc + x},
\code{self.n == 7}, \code{!self.top}, \code{sum.set(self.acc)},
\code{mux(self.top, ..)}, \code{with!(self <= \{ self.top ? .. \})} and
\code{if self.top.to\_bool()} all say what \code{get} said, the
value the last edge latched, and the \code{let} at the top of the
step that read every register into a name is gone. A drive in the
same step does not change what a read sees, as in hardware, since a
drive is deferred to the end of the step. \code{get} stays for what
needs the value as a value: a method of it, \code{self.w.get().slice}
or \code{.bit}, a typed argument, a \code{case!} or \code{select!}
over it, and a register of an enum. The lowering reads a register
named in an expression as the register, as it read \code{get}.

The unit below sums eight inputs, then holds the sum and says so on
\code{full} until a clear; the step reads its three registers only in
place, in the arithmetic, the compare, the condition and the drives
of the outputs.

\begin{figure}[htbp]
\centering
\input{inplace_timing.tex}
\caption{The window: a clear, eight inputs summed, three ignored while full, a clear and one more.}
\label{fig:x-inplace}
\end{figure}

\lstinputlisting[caption={\code{ex\_inplace.rs}. Run with \code{bazel run //lib/examples:ex\_inplace}.},label={lst:x-inplace}]{lib/examples/ex_inplace.rs}

\lstinputlisting[style=ascii,caption={What \code{ex\_inplace} printed when the build ran it: a line per cycle, then the lowering.},label={lst:x-inplace-out}]{out_inplace.txt}
